\documentclass[10pt,a4paper]{amsart}
\usepackage[margin=19mm]{geometry}
\usepackage{amsmath,amssymb,mathtools}
\usepackage[scr=rsfs,cal=euler]{mathalfa}
\usepackage{xfrac}
\usepackage[unicode,hidelinks,bookmarks=false]{hyperref}
\newcommand{\ie}{i.e.\ }
\newcommand{\cf}{cf.\ }
\newcommand{\resp}{respectively\ }
\newcommand{\Subsection}{\subsection}
\providecommand{\nobreakdash}{\nobreak\hbox{-}}
\setlength{\emergencystretch}{2em}
\multlinegap0pt
\usepackage{amssymb}
\usepackage{enumerate}
\usepackage{float}
\usepackage{mathtools}
\usepackage{xcolor}
\usepackage{bm}
\newtheorem{assumption}{Assumption}
\newtheorem{lemma}{Lemma}
\newtheorem{theorem}{Theorem}
\usepackage{cleveref}
\crefname{assumption}{Assumption}{Assumptions}
\crefname{lemma}{Lemma}{Lemmas}
\crefname{theorem}{Theorem}{Theorems}
\newcommand{\E}{\mathbb{E}}
\renewcommand{\P}{\mathbb{P}}
\renewcommand{\d}{\mathrm{d}}
\title{On the maximal displacement of subcritical branching random walk in random environment}
\author{Wenxin Fu and Wenming Hong}
\date{Source: arXiv:2605.20800v1 (CC BY 4.0)}
\begin{document}
\maketitle

\noindent\emph{Source and licence.} On the maximal displacement of subcritical branching random walk in random environment. Version arXiv:2605.20800v1 (CC BY 4.0), distributed by arXiv under the Creative Commons Attribution 4.0 International licence, http://creativecommons.org/licenses/by/4.0/, as recorded in the arXiv licence metadata for this exact version and shown on the abstract page https://arxiv.org/abs/2605.20800v1. Full licence notice: \texttt{licenses/arxiv-2605.20800-CC-BY-4.0.txt}.\par\smallskip

\noindent\textit{Editorial scope.} Counted unit: the article's lemma 2 (source0.tex lines 2186--2199). The article's complete original proof of it is delivered with it (source0.tex lines 2201--2276, one \texttt{proof} environment of 76 lines). No mathematics is written, repaired or re-derived: the statement and the proof are the article's own lines, in the article's own notation, and no retained line is substituted or re-flowed.  Retained uncounted as the source-local context the proof needs: lines 341--353; lines 367--373; lines 484--495; lines 548--562; lines 966--968; lines 2052--2066.  Nothing was omitted from any retained span, so the package carries no omission record.  No retained line contains a citation, so the wrapper carries no bibliography and no bibliography entry.  No retained line contains a figure environment, an image-inclusion command or drawing code, so no image is delivered and the package declares no asset dependency.  Editorial formatting only, in addition: an AMS \texttt{amsart} preamble that restates the article's own declarations and macros used by the retained lines, namely the environments \texttt{assumption}, \texttt{lemma}, \texttt{theorem} on separate counters, together with the standard packages those lines need.  The retained environments print in this document as Assumption 1, Lemma 1, Theorem 1 in order of appearance, whereas the article prints the same items with the numbering of its own full document.  The complete native source of the article as distributed in the arXiv e-print for this exact version is bundled unchanged as \texttt{sources/201019/source0.tex}.
	\begin{assumption}\label{subcritical}
		We assume that
		\begin{equation}
			a:=\E[X_1]\in (-\infty,0).
		\end{equation}
		and for any $\rho\geqslant0$, we have
		\begin{equation}
			\Lambda_e(\rho) = \log \E \left[ e^{\rho X} \right] <\infty.
		\end{equation}
	And define \begin{equation}\label{def of alpha}
	 \alpha:=\sup\left\{\rho:\Lambda_e(\rho)\leqslant 0\right\}.
	\end{equation}
	\end{assumption}

	\begin{assumption}\label{Jump Assumption}
		The jump distribution $\mu(\d x)$ is  supported on $\mathbb{Z}$, non-trivial, satisfies $\mu([2,\infty)) = 0$ with mean $0$. Thus, for any $\lambda>0$, denote by
		\begin{equation}
			\Lambda_s(\lambda):=\log \E_\mu\left[ e^{\lambda h}\right]<\infty,
		\end{equation}
	where under $\E_\mu$, $h$ is a random variable with distribution given by $\mu$. 
	\end{assumption}

	\begin{lemma}\label{Classify}
		Under \Cref{subcritical}and \Cref{Jump Assumption},
		the function $\varTheta(\rho)$ is finite on $(0,\alpha]$, where $\alpha:=\sup\left\{\rho:\Lambda_e(\rho)\leqslant 0\right\}$, and is a strictly increasing function on $(0,\alpha]$. Moreover, we have 
		\begin{equation}
			\lim\limits_{\rho\rightarrow 0^+} \varTheta(\rho)<0.
		\end{equation}
		And if $\alpha<\infty$, we have
		\begin{equation}
			\varTheta(\alpha)>0,
		\end{equation}
		which leads that in $ (0,\alpha] $, the exist a unique zero point of $\varTheta(\rho)$.
	\end{lemma}

	\begin{theorem}\label{Theorem of weakly cases}
		Under \Cref{subcritical} and \Cref{Jump Assumption}, assume that $\alpha>1$ and $\varTheta(1)>0$ or $\alpha<1$, by \Cref{Classify}, for both cases, there exists a unique $\rho\in(0,1)$ such that
		\begin{equation}\label{25}
			\varTheta(\rho)= \Lambda_e^\prime(\rho) + \Lambda_s(\lambda_\rho) -\lambda_\rho \Lambda_s^\prime(\lambda_\rho) = 0,
		\end{equation}
		where $\lambda_\rho$ is the unique positive solution of equation $\Lambda_e(\rho)+\rho\Lambda_s(\lambda_\rho) =0$.
		In addition, we assume that there exists some $\delta>0$ such that
		\begin{equation}\label{additional assumption for weakly}
			\E\left(e^{\rho X_1-\Lambda_e(\rho)} (\log_+ \eta_1)^{2+\delta}\right)<\infty.
		\end{equation}
		Then there exist some finite and positive constants $C_3$ and $C_4$, such that for any $x\geqslant 0$, we have
		\begin{equation}
			C_3 x^{-\frac{3}{2}} \leqslant e^{\rho\lambda_\rho x }\P(M\geqslant x ) \leqslant C_4.
		\end{equation}
	\end{theorem}

	\begin{equation}\label{def of overlineA}
		B_x:= \left(1+\sum_{k=-T_x}^{-1}e^{J_x}\Phi(\mathcal{T}_k,\overline{Y}_k,0)\right)^{-1},
	\end{equation}

	\begin{lemma}\label{Measure change Relationship 2}
		For any integer $x\geqslant 1$, for any measurable function $\phi$, we have
		\begin{equation}
			e^{\rho\lambda_\rho x}P_E\otimes P_{\overline{Y}} (e^{\rho R_x}\phi(B_x,R_x)) = Q^\rho_E\otimes P_{\overline{Y}} ( \phi(B_x,R_x) ),
		\end{equation}
		where $B_x$ is defined in \cref{def of overlineA} from branching system $(\mathcal{T},V)$.
		Especially, we have
		\begin{equation}
			e^{\rho\lambda_\rho x} P\otimes P_\xi(M\geqslant x)
			=
			e^{\rho\lambda_\rho x} P_E\otimes P_{\overline{Y}}(e^{R_x}B_x)
			=
			Q^\rho_E\otimes P_{\overline{Y}} (e^{(1-\rho)R_x}B_x).
		\end{equation}
	\end{lemma}

	\begin{lemma}\label{lemma of weakly 2}
		For any $\delta\in (0,\rho)$, there exists some constant $ C_\delta <\infty $, such that for any $x\geqslant1$, we have
		\begin{equation*}
			Q^\rho_E\otimes P_{\overline{Y}} \left(
			e^{(1-\delta)R_x} B_x
			\right) \leqslant C_\delta.
		\end{equation*}
	Then by the H\"older inequality, we have
	\begin{equation}
		Q^\rho_E\otimes P_{\overline{Y}} \left(
		e^{(1-\delta)R_x} B_x
		\right) \leqslant C_\delta^{\frac{1-\rho}{1-\delta}}<\infty.
	\end{equation}
\end{lemma}

	\begin{proof}[Proof of \Cref{lemma of weakly 2}]
	By \Cref{Measure change Relationship 2}, we have
	\begin{align}
		Q^\rho_E\otimes P_{\overline{Y}} \left(
		e^{(1-\delta)R_x} B_x
		\right) &= 
		e^{\rho\lambda_\rho x}P_E\otimes P_{\overline{Y}} \left(
		e^{(1+\rho-\delta) R_x} B_x
		\right)
		\nonumber\\
		&= e^{\rho\lambda_\rho x}
		P\otimes P_{\xi} \left( e^{(1+\rho-\delta) R_x} B_x \right).
	\end{align}
	
	Under $P\otimes \tilde{P}_\xi$, 
	for $P~a.s.$ environment $\xi$, 
	define $U:=\sup\{S_n+n\Lambda_s(\lambda_\rho);n\geqslant 0\}$.
	Note that, $\{S_n-n\Lambda_s(\lambda_\rho);n\geqslant 0\}$ is a random walk with log-Laplace exponent $ \overline{\Lambda}(\theta)=\Lambda_e(\theta)+\theta \Lambda_s(\lambda_\rho) $, which has a positive zero point at $\theta=\rho$ (under the Assumption of \Cref{Theorem of weakly cases}).
	Then, $P$ almost surely, we have $U<\infty$. 
	Moreover, we have
	\begin{equation*}
		P(U\geqslant y) \sim e^{-\rho y}.
	\end{equation*}
	Note that, if $\xi\in \{U <\lambda_\rho x\}$, we have $R_x\leqslant 0,~ \tilde{P}_\xi~a.s.$. So, we have
	\begin{align*}
		\tilde{P}_\xi(e^{(1+\rho-\delta)R_x}B_x)
		\leqslant\tilde{P}_\xi(e^{R_x}B_x),
	\end{align*}
	which implies that
	\begin{equation*}
		P\otimes \tilde{P}_\xi (e^{(1+\rho-\delta)R_x}B_x;U<\lambda x)\leqslant
		P\otimes \tilde{P}_\xi (e^{R_x}B_x).
	\end{equation*}
	
	And if $\xi\in \{U\geqslant \lambda_\rho x\}$, and the inequality $R_x\leqslant U-\lambda_1 x$ holds. Then we have
	\begin{equation*}
		\tilde{P}_\xi(e^{(1+\rho-\delta)R_x}B_x)
		\leqslant e^{(\rho-\delta)(U-\lambda_1x)}
		\tilde{P}_\xi(e^{R_x}B_x)\leqslant e^{\frac{1}{p}(U-\lambda_1x)},
	\end{equation*}
	which implies that
	\begin{equation*}
		P\otimes \tilde{P}_\xi (e^{(1+\frac{1}{p})R_x}B_x;U\geqslant\lambda x)\leqslant
		P\left( e^{\frac{1}{p}(U-\lambda_1 x)};U\geqslant \lambda_1 x \right)
	\end{equation*}
	And according to the tail probability of $U$, it is not hard to verify that as $x\rightarrow\infty$, we have
	\begin{equation*}
		P\left( e^{(\rho-\delta)(U-\lambda_1 x)}|U\geqslant \lambda_1 x \right) \sim \int_{0}^{\infty} e^{(\rho-\delta)y} e^{-\rho y} \d x <\infty.
	\end{equation*}
	Combining above relationships, there exists some finite constant $C$, such that
	\begin{equation*}
		Q_E^{\rho}\otimes P_{\overline{Y}} \left(
		e^{(1-\delta)R_x} B_x
		\right)\leqslant Q_E^{\rho} \left(
		e^{(1-\rho)R_x} B_x
		\right) +C.
	\end{equation*}
	And by H\"older inequality, we have
	\begin{equation*}
		Q_E^{\rho}\otimes P_{\overline{Y}} \left(
		e^{(1-\delta)R_x} B_x
		\right)
		\geqslant
		Q_E^{\rho}\otimes P_{\overline{Y}} \left(
		e^{(1-\rho)R_x} B_x
		\right) ^{\frac{1-\delta}{1-\rho}}
	\end{equation*}
    As long as $\delta \in (0,\rho)$, we know that the inequality $x^{\frac{1-\delta}{1-\rho}}\leqslant x+C$ implies that there exists some finite constant $C_\delta$ such that $x\leqslant C_\delta$. Therefore, we have
    \begin{equation*}
    	Q_E^{\rho}\otimes P_{\overline{Y}} \left(
    	e^{(1-\delta)R_x} B_x
    	\right) \leqslant C_\delta<\infty.
    \end{equation*}
    
    So far, the proof is finished.
    \end{proof}

\end{document}
